% ============================================================
%  III. DCPAM-CV: Classical Vision Pipeline
% ============================================================
\section{DCPAM-CV: Classical Vision Pipeline}

This section details the complete measurement pipeline
of the classical computer-vision variant.
%
Fig.~\ref{fig:pipeline} illustrates the data flow from raw
images to the final distance measurement.

% TODO: 插入 pipeline 流程图
% \begin{figure}[!t]
%   \centering
%   \includegraphics[width=\columnwidth]{pipeline}
%   \caption{DCPAM-CV measurement pipeline.}
%   \label{fig:pipeline}
% \end{figure}


% ------------------------------------------------------------
\subsection{Coordinate System Definitions}
\label{sec:coord}

We define three coordinate frames, all right-handed:

\textbf{Camera coordinate frame $C1$.}
Origin at the optical centre of camera $C_1$;
$Z$-axis aligned with the optical axis.
%
A point $P$ in this frame is denoted
$P^{C1} = (x_1, y_1, z_1)^T$.

\textbf{Camera coordinate frame $C2$.}
Origin at the optical centre of camera $C_2$;
$Z$-axis aligned with its optical axis.
%
A point in this frame is denoted
$P^{C2} = (x_2, y_2, z_2)^T$.

\textbf{Device coordinate frame $D$.}
By convention, $D \equiv C1$---the device frame coincides
with the front camera coordinate frame.
%
All final computations are performed in $D$.

The optical centres of the two cameras in their own frames are
trivially
$O_{C_1} = (0,0,0)^T$ and $O_{C_2} = (0,0,0)^T$.
%
Note that this is a definitional identity:
the intrinsic parameters (focal length, principal point,
distortion) govern the projection relationship but do not
affect the location of the optical centre in the camera's
own coordinate frame.


% ------------------------------------------------------------
\subsection{Key Points}
\label{sec:keypoints}

We define the following key points (all expressed in $C1$
unless otherwise noted):

$P^f_{\mathrm{real}}$:
the laser spot on the front receiving screen
(real image in $C1$).

$P^f_{\mathrm{virtual}}$:
the virtual-image point obtained by rotating
$P^f_{\mathrm{real}}$ about the mirror hinge
(virtual image in $C1$).

$P^{b\text{-}C2}_{\mathrm{real}}$:
the laser spot on the rear receiving screen
(real image in $C2$).

$P^{b\text{-}C2}_{\mathrm{virtual}}$:
the virtual-image point of the rear screen
(virtual image in $C2$).

$P^b_{\mathrm{virtual}}$:
$P^{b\text{-}C2}_{\mathrm{virtual}}$ transformed
into $C1$.

$P_T$:
the target point whose distance to the laser axis
is to be measured.


% ------------------------------------------------------------
\subsection{Laser Spot Centre Extraction}
\label{sec:extraction}

Let the pixel coordinates of the laser spot in
$C_1$ and $C_2$ be $(u_1, v_1)$ and $(u_2, v_2)$,
with homogeneous representations
%
\begin{equation}
  \mathbf{p}_1 =
  \begin{bmatrix} u_1 \\ v_1 \\ 1 \end{bmatrix}, \quad
  \mathbf{p}_2 =
  \begin{bmatrix} u_2 \\ v_2 \\ 1 \end{bmatrix}.
  \label{eq:pixel-coords}
\end{equation}

% TODO: 补充光斑提取的具体 CV 算法
%       （高斯拟合、质心法、Canny 边缘 + 梯度投票等，
%         参考 center.py 中的实现）


% ------------------------------------------------------------
\subsection{Pixel-to-Camera Coordinate Transformation}
\label{sec:pixel2cam}

Let $K_1$ and $K_2$ be the intrinsic matrices of cameras
$C_1$ and $C_2$, and let $Z_s$ be the known distance from
the camera optical centre to the receiving screen along the
optical axis.
%
The back-projection from pixel coordinates to camera
coordinates is
%
\begin{equation}
  \boxed{
    P^f_{\mathrm{real}} = Z_s \, K_1^{-1}\, \mathbf{p}_1, \quad
    P^{b\text{-}C2}_{\mathrm{real}}
      = Z_s \, K_2^{-1}\, \mathbf{p}_2
  }
  \label{eq:backproj}
\end{equation}

\textbf{Dependencies.}
Camera calibration provides $K_1$, $K_2$;
the tool geometry provides $Z_s$.

\textbf{Error sources.}
The accuracy of $Z_s$ depends on the mechanical precision
of the probe fixture.


% ------------------------------------------------------------
\subsection{Real-to-Virtual Point Transformation}
\label{sec:real2virtual}

The receiving screen is mounted on a mirror that can rotate
about a hinge axis.
%
To reconstruct the ``unfolded'' laser path,
each real image point must be rotated to its virtual position.

In frame $C1$, let the rotation centre be
$R_{O1} = (x_{O1}, 0, z_{O1})^T$,
and let $\theta$ be the rotation angle
(counter-clockwise from $Z$ towards $X$ in a right-handed
frame).
%
Then
%
\begin{equation}
  \boxed{P^f_{\mathrm{virtual}} = H_1 \cdot P^f_{\mathrm{real}}}
  \label{eq:H1}
\end{equation}
%
where $H_1$ is a $4\times4$ rigid-body transformation matrix
combining rotation and translation:
%
\begin{equation}
  H_1 =
  \begin{bmatrix}
    \cos\theta  & 0 & \sin\theta
      & x_{O1}(1{-}\cos\theta) - z_{O1}\sin\theta \\
    0           & 1 & 0           & 0 \\
    -\sin\theta & 0 & \cos\theta
      & z_{O1}(1{-}\cos\theta) + x_{O1}\sin\theta \\
    0           & 0 & 0           & 1
  \end{bmatrix}
  \label{eq:H1-matrix}
\end{equation}

The transformation represents a rotation about an axis
parallel to $Y$ that passes through $R_{O1}$.

Similarly, the rear real-image point is rotated in $C2$
about the centre $R_{O2} = (x_{O2}, 0, z_{O2})^T$
by angle $2\theta$:
%
\begin{equation}
  P^{b\text{-}C2}_{\mathrm{virtual}}
    = H_2 \cdot P^{b\text{-}C2}_{\mathrm{real}}
  \label{eq:H2}
\end{equation}

\textbf{Dependencies.}
Geometric dimensions $R_{O1}$, $R_{O2}$, and the
rotation angle $\theta$.

\textbf{Ideal-case simplification.}
Under ideal assembly,
$z_{O1} = z_{O2} = Z_s$.

% TODO: 补充旋转推导过程的详细数学证明（可放入附录）


% ------------------------------------------------------------
\subsection{$C2$-to-$C1$ Coordinate Transformation}
\label{sec:c2toc1}

The virtual-image point $P^{b\text{-}C2}_{\mathrm{virtual}}$
is expressed in frame $C2$ and must be converted to $C1$:
%
\begin{equation}
  \boxed{
    P^b_{\mathrm{virtual}}
      = T_{C2 \to C1}
        \cdot P^{b\text{-}C2}_{\mathrm{virtual}}
  }
  \label{eq:T-c2c1}
\end{equation}

The transformation matrix is computed from the extrinsic
parameters of the two cameras with respect to a common world
frame $W$:
%
\begin{equation}
  T_{C2 \to C1} = E_1 \cdot E_2^{-1}
  \label{eq:T-from-ext}
\end{equation}
%
where $E_1$, $E_2$ are the $4\times4$ extrinsic matrices
(world-to-camera) of $C_1$ and $C_2$ respectively.

Equivalently, if the relative pose is parameterised as
rotation $R_{21}$ and translation $t_{21}$
(from $C1$ to $C2$), then
%
\begin{align}
  R_{12} &= R_{21}^T  \label{eq:R12}\\
  t_{12} &= -R_{21}^T \, t_{21}  \label{eq:t12}
\end{align}
%
and $P^{C1} = R_{12}\, P^{C2} + t_{12}$.

\textbf{Dependencies.}
Stereo camera calibration (extrinsic parameters).


% ------------------------------------------------------------
\subsection{Point-to-Axis Distance Computation}
\label{sec:distance}

With $P^f_{\mathrm{virtual}}$, $P^b_{\mathrm{virtual}}$,
and $P_T$ all expressed in $C1$,
the perpendicular distance from $P_T$ to the laser axis
(the line through
$P^f_{\mathrm{virtual}}$ and $P^b_{\mathrm{virtual}}$)
is given by the classical cross-product formula:
%
\begin{equation}
  \boxed{
    H = \frac{
      \left\| \vec{L} \times \vec{w} \right\|
    }{
      \left\| \vec{L} \right\|
    }
  }
  \label{eq:distance}
\end{equation}
%
where
%
\begin{align}
  \vec{L} &= P^f_{\mathrm{virtual}} - P^b_{\mathrm{virtual}}
  \label{eq:L-vec} \\
  \vec{w} &= P_T - P^b_{\mathrm{virtual}}
  \label{eq:w-vec}
\end{align}

The vector $\vec{L}$ is the direction vector of the
reconstructed laser axis in frame $C1$.

\textbf{Dependencies.}
The target point $P_T$ is determined by the tool-bar
geometry and the measurement configuration.
